\section{Discussions}
\label{sec:discussions}

Currently, \sys only supports the forward pass for attention computation.
To extend \sys and apply it to training would require developing customizable backward attention kernel templates, which we plan to explore in future work.
Regarding the generality of \sys, our approach decouples computation from tile scheduling, allowing for diverse tiling strategies such as FlashDecoding~\cite{flashdecodingpp} and Lean Attention~\cite{sanovar2024lean} by delegating scheduling to a runtime-scheduler rather than embedding it within the attention template.
This design generalizes scheduling algorithms, as detailed in Algorithm \ref{alg:load-balancing}, for load-balancing and wave-quantization reduction while targeting optimal GPU performance across architectures (e.g., FlashAttention2~\cite{dao2023flashattention2} for Turing/Ampere/Ada and FlashAttention3~\cite{jay2024flashattention3} for Hopper).
Although frameworks like Triton~\cite{triton} offer GPU-agnostic interfaces, they often lag in adopting new hardware features; our template design space, expressed as $f_\textrm{epilogue}(\textrm{scan}(f_\textrm{logits}(f_{q}(Q)\cdot f_{k}(K)))\cdot f_{v}(V))$, covers most attention functions, including recent variants such as Multi-head Latent Attention (MLA)~\cite{deepseekv2} and the intra-attention component of Linear Attention~\cite{yang2024gatedlinearattentiontransformers}.
